 In the above diagram, we can see that there is a connection between the legs of the same size $R$. This is consistent with matrix multiplication where there is a sum over indices with the same dimension. Finally, since the resulting diagram has two free legs, it represents a matrix. Here, the final object is a matrix which can be represented as  a single node, demonstrating how nodes can be merged in tensor networks:
$$
\Ab\Bb =\Mb \Leftrightarrow 
\begin{tikzpicture}[baseline=\baseline]
   \tikzset{tensor/.style = {minimum size = 0.5cm,shape = circle,thick,draw=black,inner sep = 0pt}, edge/.style = {thick,line width=.4mm},every loop/.style={}}
    \def\y{0}
    \def\x{0}
    \node[tensor,fill=red!50!white] (A) at (\x,\y){\scalebox{0.85}{$\Ab$}};
    \node[tensor,fill=blue!50!white] (B) at (\x+1,\y){\scalebox{0.85}{$\Bb$}};
    \draw[edge] (A) -- (B) node [above,midway] {\scalebox{0.7}{\dimcolor{$R$}}};
    \draw[edge] (\x-0.25,\y) -- (\x-0.75,\y)  node [midway, above] {\scalebox{0.7}{\dimcolor{$d_1$}}};
    \draw[edge] (\x+1.25,\y) -- (\x+1.75,\y)  node [midway,above] {\scalebox{0.7}{\dimcolor{$d_2$}}};
    \node[draw=none] () at (\x+2,\y) {\scalebox{1}{\textcolor{black}{$=$}}};
    \node[tensor,fill=blue!30!red!20!white] (M) at (\x+3,\y){\scalebox{0.85}{$\Mb$}};
    \draw[edge] (M) -- (\x+3.75,\y)  node [midway,above] {\scalebox{0.7}{\dimcolor{$d_2$}}};
    \draw[edge] (M) -- (\x+2.25,\y)  node [midway,above] {\scalebox{0.7}{\dimcolor{$d_1$}}};
    \end{tikzpicture}.
$$
Contraction between two matrices can directly be extended to contractions between a matrix and a vector, a tensor and a matrix, a tensor a matrix and a vector products, etc. Here are some examples of tensor networks with their corresponding mathematical expressions: 
\begin{align}\label{eq:matvec}
\Ab\in\RR^{d_1\times R},\vb\in\RR^{R} \ : \ \
\begin{tikzpicture}[baseline=\baseline]
   \tikzset{tensor/.style = {minimum size = 0.5cm,shape = circle,thick,draw=black,inner sep = 0pt}, edge/.style = {thick,line width=.4mm},every loop/.style={}}
    \def\y{0}
    \def\x{0}
    \node[tensor,fill=red!50!white] (A) at (\x,\y){\scalebox{0.85}{$\Ab$}};
    \node[tensor,fill=blue!10!white] (v) at (\x+1,\y){\scalebox{0.85}{$\vb$}};
    \draw[edge] (A) -- (v) node [above,midway] {\scalebox{0.7}{\dimcolor{$R$}}};
    \draw[edge] (\x-0.25,\y) -- (\x-0.75,\y)  node [pos=1.3]{\scalebox{0.7}{\indcolor{$i$}}};
    \end{tikzpicture}
=\sum_{r=1}^R\Ab_{ir}\vb_r,
\end{align}
\begin{align}\label{eq:tenmat}
\At\in\RR^{d_1\times R\times d_2},\Ab\in\RR^{R\times d_3}\ :\ \
\begin{tikzpicture}[baseline = \baseline]
    \tikzset{tensor/.style = {minimum size = 0.5cm,shape = circle,thick,draw=black,inner sep = 0pt}, edge/.style = {thick,line width=.4mm},every loop/.style={}}
    \def\y{0}
    \def\x{0}
    \node[tensor,fill=green!50!red!50!white] (At) at (\x,\y){\scalebox{0.85}{$\At$}};
    \node[tensor,fill=red!50!white] (A) at (\x+1,\y){\scalebox{0.85}{$\Ab$}};
    \draw[edge] (At) -- (A) node [above,midway] {\scalebox{0.7}{\dimcolor{$R$}}};
    \draw[edge] (\x-0.25,\y) -- (\x-0.75,\y)  node [pos=1.3] {\scalebox{0.7}{\indcolor{$i$}}};
    \draw[edge] (\x,\y-0.25) -- (\x,\y-0.75)  node [pos=1.3] {\scalebox{0.7}{\indcolor{$j$}}};
    \draw[edge] (\x+1.25,\y) -- (\x+1.75,\y)  node [pos=1.3] {\scalebox{0.7}{\indcolor{$k$}}};
\end{tikzpicture} 
= \sum_{r=1}^R\At_{ijr}\Ab_{rk},
\end{align}
\begin{align}\label{eq:tenmatvec}
\At\in\RR^{d_1\times R\times d_2},\Ab\in\RR^{R\times d_3}, \vb\in\RR^{d_3}&:
   \begin{tikzpicture}[baseline = \baseline]
    \tikzset{tensor/.style = {minimum size = 0.5cm,shape = circle,thick,draw=black,inner sep = 0pt}, edge/.style = {thick,line width=.4mm},every loop/.style={}}
    \def\y{0}
    \def\x{0}
    \node[tensor,fill=green!50!red!50!white] (At) at (\x,\y){\scalebox{0.85}{$\At$}};
    \node[tensor,fill=red!50!white] (A) at (\x+1,\y){\scalebox{0.85}{$\Ab$}};
    \node[tensor,fill=blue!40!red!60!white] (a) at (\x+2,\y){\scalebox{0.85}{$\vb$}};
    \draw[edge] (At) -- (A) node [above,midway] {\scalebox{0.7}{\dimcolor{$R$}}};
    \draw[edge] (A) -- (a) node [above,midway] {\scalebox{0.7}{\dimcolor{$d_3$}}};
    \draw[edge] (\x-0.25,\y) -- (\x-0.75,\y)  node [pos=1.3] {\scalebox{0.7}{\indcolor{$i$}}};
    \draw[edge] (\x,\y-0.25) -- (\x,\y-0.75)  node [pos=1.3] {\scalebox{0.7}{\indcolor{$j$}}};
    \end{tikzpicture} 
= \sum_{r=1}^R\sum_{k=1}^{d_3}\At_{ijr}\Ab_{rk}\ab_{k}.
\end{align}
As we can see in all diagrams above, edges between two nodes represent summations. They also indicate that two tensors share dimensions of the same size on the corresponding mode. 
As explained before, the number of free legs in the tensor network corresponds to the order of the tensor it represents. For the previous examples, we have

$$
\begin{tikzpicture}[baseline=\baseline]
   \tikzset{tensor/.style = {minimum size = 0.5cm,shape = circle,thick,draw=black,inner sep = 0pt}, edge/.style = {thick,line width=.4mm},every loop/.style={}}
    \def\y{0}
    \def\x{0}
    \node[tensor,fill=red!50!white] (A) at (\x,\y){\scalebox{0.85}{$\Ab$}};
    \node[tensor,fill=blue!10!white] (v) at (\x+1,\y){\scalebox{0.85}{$\vb$}};
    \draw[edge] (A) -- (v) node [above,midway] {\scalebox{0.7}{\dimcolor{$R$}}};
    \draw[edge] (\x-0.25,\y) -- (\x-0.75,\y)  node [midway, above] {\scalebox{0.7}{\dimcolor{$d_1$}}};
    \end{tikzpicture} 
   \in \mathbb R^{d_1},\ \ 
\begin{tikzpicture}[baseline = \baseline]
    \tikzset{tensor/.style = {minimum size = 0.5cm,shape = circle,thick,draw=black,inner sep = 0pt}, edge/.style = {thick,line width=.4mm},every loop/.style={}}
    \def\y{0}
    \def\x{0}
    \node[tensor,fill=green!50!red!50!white] (At) at (\x,\y){\scalebox{0.85}{$\At$}};
    \node[tensor,fill=red!50!white] (A) at (\x+1,\y){\scalebox{0.85}{$\Ab$}};
    \draw[edge] (At) -- (A) node [above,midway] {\scalebox{0.7}{\dimcolor{$R$}}};
    \draw[edge] (\x-0.25,\y) -- (\x-0.75,\y)  node [midway,above] {\scalebox{0.7}{\dimcolor{$d_1$}}};
    \draw[edge] (\x,\y-0.25) -- (\x,\y-0.75)  node [midway,left] {\scalebox{0.7}{\dimcolor{$d_2$}}};
    \draw[edge] (\x+1.25,\y) -- (\x+1.75,\y)  node [midway,above] {\scalebox{0.7}{\dimcolor{$d_3$}}};
\end{tikzpicture} 
\in \mathbb R^{d_1\times d_2 \times d_3},\ \ \text{and}\ \ 
    \begin{tikzpicture}[baseline = \baseline]
    \tikzset{tensor/.style = {minimum size = 0.5cm,shape = circle,thick,draw=black,inner sep = 0pt}, edge/.style = {thick,line width=.4mm},every loop/.style={}}
    \def\y{0}
    \def\x{0}
    \node[tensor,fill=green!50!red!50!white] (At) at (\x,\y){\scalebox{0.85}{$\At$}};
    \node[tensor,fill=red!50!white] (A) at (\x+1,\y){\scalebox{0.85}{$\Ab$}};
    \node[tensor,fill=blue!40!red!60!white] (a) at (\x+2,\y){\scalebox{0.85}{$\vb$}};
    \draw[edge] (At) -- (A) node [above,midway] {\scalebox{0.7}{\dimcolor{$R$}}};
    \draw[edge] (A) -- (a) node [above,midway] {\scalebox{0.7}{\dimcolor{$d_3$}}};
    \draw[edge] (\x-0.25,\y) -- (\x-0.75,\y)  node [midway,above] {\scalebox{0.7}{\dimcolor{$d_1$}}};
    \draw[edge] (\x,\y-0.25) -- (\x,\y-0.75)  node [midway,left] {\scalebox{0.7}{\dimcolor{$d_2$}}};
    \end{tikzpicture} 
\in \mathbb R^{d_1\times d_2}.
$$
 
\paragraph{From mathematical expressions to tensor networks, and back} Tensor networks are simple to work with because there are no strict rules for representing legs and nodes. 
In tensor network diagrams, nodes and edges can be positioned arbitrarily in the plane, only the graph structure matters. 
For example, when translating matrix multiplication into  a tensor network diagram, the key feature is to ensure the dimension of the connected legs are consistent, e.g., for $\Ab\in\RR^{d_1\times R}, \Bb\in\RR^{R\times d_2}$  all diagrams below illustrate the \emph{same} matrix multiplication:
\begin{align*} 
\Ab\Bb
=
\begin{tikzpicture}[baseline=\baseline]
   \tikzset{tensor/.style = {minimum size = 0.5cm,shape = circle,thick,draw=black,inner sep = 0pt}, edge/.style = {thick,line width=.4mm},every loop/.style={}}
    \def\y{0}
    \def\x{0}
    \node[tensor,fill=red!50!white] (A) at (\x,\y){\scalebox{0.85}{$\Ab$}};
    \node[tensor,fill=blue!50!white] (B) at (\x+1,\y){\scalebox{0.85}{$\Bb$}};
    \draw[edge] (A) -- (B) node [midway,above] {\scalebox{0.7}{\dimcolor{$R$}}};;
    \draw[edge] (\x-0.25,\y) -- (\x-0.75,\y) node [midway,above] {\scalebox{0.7}{\dimcolor{$d_1$}}};;
    \draw[edge] (\x+1.25,\y) -- (\x+1.75,\y) node [midway,above] {\scalebox{0.7}{\dimcolor{$d_2$}}};;
    \node[draw=none] () at (\x+1.95,\y) {$=$};
    \node[tensor,fill=red!50!white] (A1) at (\x+2.85,\y){\scalebox{0.85}{$\Ab$}};
    \node[tensor,fill=blue!50!white] (B1) at (\x+3.85,\y){\scalebox{0.85}{$\Bb$}};
    \draw[edge] (A1) -- (B1) node [midway,above] {\scalebox{0.7}{\dimcolor{$R$}}};;
    \draw[edge] (A1) to [out=60,in=60, distance=0.5cm] (\x+2.2,\y);
    \draw[edge] (B1) to [out=60,in=60, distance=0.5cm] (\x+4.4,\y-0.5);
    \node[draw=none] () at (\x+2.4,\y+0.65) {\scalebox{0.7}{\dimcolor{$d_1$}}};
    \node[draw=none] () at (\x+4,\y+0.65) {\scalebox{0.7}{\dimcolor{$d_2$}}};
    \node[draw=none] () at (\x+4.7,\y) {$=$};
    \node[tensor,fill=blue!50!white] (A2) at (\x+5.3,\y){\scalebox{0.85}{$\Bb$}};
    \node[tensor,fill=red!50!white] (B2) at (\x+6.3,\y){\scalebox{0.85}{$\Ab$}};
    \draw[edge] (A2) -- (B2) node [midway,above] {\scalebox{0.7}{\dimcolor{$R$}}};;
    \draw[edge] (A2) -- (\x+5.3,\y+0.75) node [midway,left] {\scalebox{0.7}{\dimcolor{$d_2$}}};
     \draw[edge] (B2) -- (\x+6.3,\y-0.75) node [midway,right] {\scalebox{0.7}{\dimcolor{$d_1$}}};   
     \node[draw=none] () at (\x+7.1,\y) {$=$};
    \node[tensor,fill=red!50!white] (A3) at (\x+7.6,\y+0.5){\scalebox{0.85}{$\Ab$}};
    \node[tensor,fill=blue!50!white] (B3) at (\x+7.6,\y-0.5){\scalebox{0.85}{$\Bb$}};
    \draw[edge] (A3) -- (B3) node [midway,right] {\scalebox{0.7}{\dimcolor{$R$}}};
    \draw[edge] (A3) -- (\x+7.6,\y+1.25) node [midway,left] {\scalebox{0.7}{\dimcolor{$d_1$}}};
    \draw[edge] (B3) -- (\x+7.6,\y-1.25) node [midway,left] {\scalebox{0.7}{\dimcolor{$d_2$}}};
    \end{tikzpicture}.
\end{align*}
Note that because of the sizes of $\Ab\in\RR^{d_1\times R}$ and $\Bb\in\RR^{R\times d_2}$, there is only one way to do the contraction between the two; hence, there would be no ambiguity in the diagrams above even if we omitted the dimensions on the edges. 

